\documentclass{tufte-handout}

\usepackage{style}

\begin{document}

\begin{question}

    Question 8 from section B

Let \( R \) be the ring \( \Z[\sqrt{-13}] \)

\qpart

\[ A = \{ s + t\sqrt{-13} : s,t \in \Z, s - t \equiv 0\pmod{7} \} \]

As \( s = t\pmod{7} \) there eisits an integer \( k \) such that,
\[ s = t + 7k \]
Hence,
\begin{align*}
s + t\sqrt{-13} &= (t + 7k) + t\sqrt{-13} \\[8pt]
&= t(1 + \sqrt{-13}) + 7k\\[8pt]
\stext{So}
A \langle 7, 1 + \sqrt{-13} \rangle
\end{align*} 

By \textup{Definition 2.1 HB Chapter 17 p94},

To Show \( A \) is ideal, it must be closed under multiplication by any \( r \in R \).
So we need the product of the generator and \( \sqrt{-13} \) to stay in \( A \).

\begin{align*}
\sqrt{-13}(1 + \sqrt{-13}) &= \sqrt{-13} -13\\[8pt]
\stext{since} -13 &\equiv 1\pmod{7} \\[8pt]
-13 &= 1 -14\\[8pt]
\sqrt{-13} -13 &= (\sqrt{-13} + 1) - 14\\[8pt]
&= (\sqrt{-13} + 1) + 7(-2)
\end{align*}

Since \( (\sqrt{-13} + 1) \in A \) and \( 14 = 2(7) \in A \) it follows that \( A \) 
is closed under multiplication by any element of \( R \) and hence is an ideal.

\vspace{3cm}

\qpart

\[ B = \{ u + v\sqrt{-13} : u,v \in \Z, u + v \equiv 0\pmod{7} \} \]

As \( u = -v\pmod{7} \) there eisits an integer \( k \) such that, 
\[ u = -v + 7k \] 

Hence, 
\begin{align*} 
    u + v\sqrt{-13} &= (-v + 7k) + v\sqrt{-13} \\[8pt] 
    &= v(-1 + \sqrt{-13}) + 7k\\[8pt] 
    \stext{So} 
    B \langle 7, -1 + \sqrt{-13} \rangle
\stext{because \( -1 \) is a unit in \( \Z[\sqrt{-13}] \), therefore we can multiply the second generator by \( -1 \) and it will still generate the same ideal.}
    B \langle 7, 1 - \sqrt{-13} \rangle
\end{align*}

Hence,

\begin{align*}
A & \langle 7, 1 + \sqrt{-13} \rangle \\[8pt]
\stext{and}
B & \langle 7, 1 - \sqrt{-13} \rangle\\[8pt]
\stext{By \textup{Lemma 3.9 HB Chapter 17 p97}, the product of two ideals \(A\) and \( B \)}
AB & \langle 7\cdot 7, 7(1 - \sqrt{-13}), 7(1 + \sqrt{-13}), (1 + \sqrt{-13})(1 - \sqrt{-13}) \rangle \\[8pt] 
& \langle 49, 7(1 - \sqrt{-13}), 7(1 + \sqrt{-13}), 14 \rangle\\[8pt]
\stext{Since the \( \hcf(14, 49) = 7 \) and \( 7(1 - \sqrt{-13}) \) and \( 7(1 + \sqrt{-13}) \) are multiples of \( 7 \) and,}
7 &= 49 -3(14) \in AB\\[8pt]
7 &\subseteq AB\\[8pt]
\stext{hence}
AB & \langle 7 \rangle
\end{align*}



\end{question}

\end{document}